\ifdefined\gkver\else\def\gkver{student}\fi
\documentclass[\gkver]{gaokaozhenti}
\begin{document}

\chapter{2000年广东}
%% paperType: 理综物理部分

\begin{enumerate}
\item
%% number: 3
%% paperName: 2000年广东高考第 3 题 4 分
%% typeId: 1
%% type: 单选题
%% chapter: 力物体的平衡
%% point: 弹力
%% method: 无
%% score: 4
%% degree: 550
%% duplicateId: 0
%% body:
$S_{1}$ 和 $S_{2}$ 分别表示劲度系数为 $k_{1}$ 和 $k_{2}$ 的两根弹簧，$k_{1} > k_{2}$。a 和 b 表示质量分别为 $m_{a}$ 和 $m_{b}$ 的两个小物体，$m_{a} > m_{b}$，将弹簧与物块按图\ref{2000广东01}所示方式悬挂起来，现要求两根弹簧的总长度最大，则应使 \xzanswer{D}
\onepicture[num=1, label=2000广东01, w=2.2cm]{\includesvg{figs/03}}
\fourchoices[answer=D]
{$S_{1}$ 在上，a 在上}
{$S_{1}$ 在上，b 在上}
{$S_{2}$ 在上，a 在上}
{$S_{2}$ 在上，b 在上}

%% answer: D
%% memo:
\memoanswer{本题原卷及材料中未提供详解，待补充。}

\item
%% number: 5
%% paperName: 2000年广东高考第 5 题 4 分
%% typeId: 2
%% type: 多选题
%% chapter: 气体性质
%% point: 液柱移动定性分析
%% method: 无
%% score: 4
%% degree: 550
%% duplicateId: 0
%% body:
如图\ref{2000广东02}所示，一气缸竖直倒放，气缸内有一质量不可忽略的活塞，将一定量的理想气体封在气缸内，活塞与气缸壁无摩擦，气体处于平衡状态，现保持温度不变把气缸稍微倾斜一点，在达到平衡后，与原来相比，则 \xzanswer{AD}
\onepicture[num=2, label=2000广东02, w=3.4cm]{\includesvg{figs/05}}
\fourchoices[answer=AD]
{气体的压强变大}
{气体的压强变小}
{气体的体积变大}
{气体的体积变小}

%% answer: AD
%% memo:
\memoanswer{本题原卷及材料中未提供详解，待补充。}

\item
%% number: 7
%% paperName: 2000年广东高考第 7 题 4 分
%% typeId: 2
%% type: 多选题
%% chapter: 机械振动
%% point: 振动图象
%% method: 无
%% score: 4
%% degree: 850
%% duplicateId: 0
%% body:
一列简谐横波沿 $x$ 轴正方向传播，在 $x$ 轴上相距 $2\Ucm$ 的 P 点和 Q 点的振动图线均如图\ref{2000广东03}所示，由此可以确定这列波的 \xzanswer{AC}
\onepicture[num=3, label=2000广东03, w=5.5cm]{figs/07.png}
\fourchoices[answer=AC]
{振幅}
{波长}
{频率}
{波速}

%% answer: AC
%% memo:
\memoanswer{本题原卷及材料中未提供详解，待补充。}

\item
%% number: 9
%% paperName: 2000年广东高考第 9 题 4 分
%% typeId: 2
%% type: 多选题
%% chapter: 电路
%% point: 动态分析
%% method: 无
%% score: 4
%% degree: 450
%% duplicateId: 0
%% body:
图\ref{2000广东04}中 A 为理想电流表，$V_{1}$ 和 $V_{2}$ 为理想电压表，$R_{1}$ 为定值电阻，$R_{2}$ 为可变电阻，电池 $E$ 内阻不计，则 \xzanswer{BCD}
\onepicture[num=4, label=2000广东04, w=4.8cm]{\includesvg{figs/09}}
\fourchoices[answer=BCD]
{$R_{2}$ 不变时，$V_{2}$ 读数与 A 读数之比等于 $R_{1}$}
{$R_{2}$ 不变时，$V_{1}$ 读数与 A 读数之比等于 $R_{1}$}
{$R_{2}$ 改变一定量时，$V_{2}$ 读数的变化量与 A 读数的变化量之比的绝对值等于 $R_{1}$}
{$R_{2}$ 改变一定量时，$V_{1}$ 读数的变化量与 A 读数的变化量之比的绝对值等于 $R_{1}$}

%% answer: BCD
%% memo:
\memoanswer{本题原卷及材料中未提供详解，待补充。}

\item
%% number: 19
%% paperName: 2000年广东高考第 19 题 13 分
%% typeId: 6
%% type: 计算题
%% chapter: 功和能
%% point: 动能定理
%% method: 无
%% score: 13
%% degree: 350
%% duplicateId: 0
%% body:
面积很大的水池，水深为 $H$，水面上浮着一正方体木块，木块边长为 $a$，密度为水的 $\frac{1}{2}$，质量为 $m$，开始时，木块静止，有一半没入水中，如图\ref{2000广东05}所示，现用力 $F$ 将木块缓慢地压到池底，不计摩擦，求：
\begin{enumerate}
	\item
	从木块刚好完全没入水中到停在池底的过程中，池水势能的改变量。

	\item
	从开始到木块刚好完全没入水的过程中，力 $F$ 所做的功。
\end{enumerate}
\onepicture[num=5, label=2000广东05, w=4.2cm, align=right]{\includesvg{figs/19}}
%% answer:
\jdanswer{
\begin{enumerate}
	\item
	$\Delta E = 2mg(H - a)$

	\item
	$W = \frac{1}{4}mga$
\end{enumerate}
}
%% memo:
\memoanswer{本题原卷及材料中未提供详解，待补充。}

\item
%% number: 21
%% paperName: 2000年广东高考第 21 题 13 分
%% typeId: 6
%% type: 计算题
%% chapter: 磁场
%% point: 带电粒子在复合场中的运动
%% method: 无
%% score: 13
%% degree: 350
%% duplicateId: 0
%% body:
图\ref{2000广东06}为一种可用于测量电子电量 $e$ 与质量 $m$ 比值 $\frac{e}{m}$ 的阴极射线管，管内处于真空状态，图中 L 是灯丝，当接上电源时可发出电子。A 是中央有小圆孔的金属板，当 L 和 A 间加上电压时（其电压值比灯丝电压大很多），电子将被加速并沿图中虚直线所示的路径到达荧光屏 S 上的 O 点，发出荧光，$P_{1}$、$P_{2}$ 为两块平行于虚直线的金属板，已知两板间距为 $d$，在虚线所示的圆形区域内可施加一匀强磁场，已知其磁感强度为 $B$，方向垂直纸面向外，a、$b_{1}$、$b_{2}$、$c_{1}$、$c_{2}$ 都是固定在管壳上的金属引线，$E_{1}$、$E_{2}$、$E_{3}$ 是三个电压可调并可读出其电压值的直流电源。
\begin{enumerate}
	\item
	试在图中画出三个电源与阴极射线管的有关引线的连线。

	\item
	导出计算 $\frac{e}{m}$ 的表达式，要求用应测物理量及题给已知量表示。
\end{enumerate}
\onepicture[num=6, label=2000广东06, w=6.5cm, align=right]{figs/21.png}
%% answer:
\jdanswer{
\begin{enumerate}
	\item
	各电源的连线如图所示。

	\item
	$\frac{U_{3}^{2}}{2U_{2}B^{2}d^{2}}$
\end{enumerate}
\onepicture[w=6cm]{\includesvg{figs/21_an}}
}
%% memo:
\memoanswer{本题原卷及材料中未提供详解，待补充。}

\end{enumerate}

\end{document}
